\section{Results}
\label{sec:results}

\subsection{Main Ablation Matrix}

\begin{table}[htbp]
\centering
\input{tables/component_necessity_ablation}
\caption{Seven-configuration component ablation in the custom scaled-governance stress workload (artifact means). Reward values in this table are workload-specific and not directly comparable to constrained-baseline or efficiency-suite reward scales.}
\label{tab:component-necessity}
\end{table}

\begin{figure}[H]
\centering
\includegraphics[width=0.95\linewidth]{figures/money_plot_ablation.pdf}
\caption{Four-configuration ablation subset plot (Full, Classical PPO, No Q-Gate, No Shield): reward, safety violations, target MAE, and fallback activation.}
\label{fig:money-plot}
\end{figure}

Figure~\ref{fig:money-plot} only plots four lines so the page stays readable. Table~\ref{tab:component-necessity} carries the full seven rows, including No Fallback, Heuristic, and Random---that last row is where the negative result lives.

The ablation table is blunt arithmetic. Dropping the quantum gate drags reward from $-1264.90$ to $-1294.06$ and pushes fallback from $0.000$ to $0.666$. Dropping the shield (or fallback in this workload) blows tracking MAE from $3.372$ to $4.063$. Classical PPO sits at $-1284.05$ on reward versus $-1264.90$ for the full stack and shows higher latency in the artifact logs.

Random is the awkward line: $-1232.33$ reward, which beats the full system on that scalar. We are not hiding it as a typo. Random can surf reward-shaping quirks without honoring control intent; the gated stack often buys safety and tracking by giving away return. That is exactly why we print violations and MAE next to reward instead of crowning a single leaderboard winner.

\subsection{Governance Controller Comparison}

\begin{table}[htbp]
\centering
\input{tables/constrained_baseline_comparison}
\caption{Governance controller comparison under the 50-seed matched protocol across scaled and DeFi environments: Random, Heuristic, CPO-style, P3O-style, PPO-Lagrangian-style, and Robust QGate+Shield. Reward values in this table are protocol-specific and not numerically comparable to ablation or efficiency-suite reward scales.}
\label{tab:constrained-baselines}
\end{table}

Fifty seeds, matched protocol, six controllers---nobody sweeps every column in both environments. Scaled governance still shows heuristic and robust gated controllers trading blows with the CPO/P3O/Lagrangian-style entries on reward without blowing up violations. In the DeFi table P3O-style wins mean reward, while the robust gated row owns the lowest target MAE, i.e., the usual reward--tracking tug of war.

The headline we want from that split is not ``we topped P3O everywhere.'' It is parameter efficiency with similar safety posture: the separate efficiency suite puts similar safety numbers next to a $166.89\times$ smaller policy parameterization. You can also read the constrained block as a cautionary tale---lower final safety cost does not automatically come with higher return once budgets bite, and several seeds move backward on return when safety is enforced hard.

\subsection{Primary Standard Benchmark Evaluation}

\begin{table}[htbp]
\centering
\footnotesize
\setlength{\tabcolsep}{3pt}
\begin{tabular}{p{0.27\linewidth}p{0.16\linewidth}p{0.47\linewidth}}
\toprule
Benchmark & Status & Note \\
\midrule
SPG1-v0 & 30 seeds done & Return mean$\pm$std: CPO -0.42$\pm$0.89, P3O -0.14$\pm$1.01, PPOLag -0.16$\pm$1.15; Cost mean$\pm$std: 65.77$\pm$70.45, 65.69$\pm$69.28, 58.08$\pm$60.34 \\
SHCV-v1 & 30 seeds done & Return mean$\pm$std: CPO -622.28$\pm$51.90, P3O -659.60$\pm$54.03, PPOLag -647.55$\pm$48.17; Cost mean$\pm$std: CPO 0.00$\pm$0.00, P3O 0.00$\pm$0.00, PPOLag 0.01$\pm$0.06 \\
\bottomrule
\end{tabular}

\caption{Completed 30-seed constrained benchmark runs on this host (SafetyPointGoal1-v0, SafetyHalfCheetahVelocity-v1).}
\label{tab:standard-transfer-status}
\end{table}

Table~\ref{tab:standard-transfer-status} is the primary external-validity result for this paper. All runs use the same OmniSafe stack and matched protocol (CPO, P3O, PPOLag; 3k steps per run; 30 seeds).

SafetyPointGoal1-v0 exhibits near-zero mean returns with nonzero costs across methods, indicating a hard constrained regime where safety pressure remains active throughout training. Under this same protocol, PPOLag reduces constraint cost by about 11.6\% relative to P3O while accepting about a 14.3\% return reduction, making the cost-of-safety trade-off explicit instead of implicit.

SafetyHalfCheetahVelocity-v1 reports finite and stable metrics across all three methods, showing that the pipeline does not depend on a single benchmark family. Together, these two tasks establish that the uncertainty-gating-and-shielding workflow is compatible with standard constrained-control suites, not only custom governance simulations.

\subsection{Stress Sweep}

\begin{figure}[H]
\centering
\includegraphics[width=0.88\linewidth]{figures/stress_sweep_top_tier.pdf}
\caption{Reward degradation under increasing adversarial intensity.}
\label{fig:stress-top-tier}
\end{figure}

In this stress sweep, the robust controller overtakes classical PPO in the higher-intensity regime (approximately $C\ge 1.4$ in the associated complexity summary), while deterministic heuristic control remains competitive across much of the range. The key pattern is regime-dependent ranking, not a single globally dominant method: deterministic heuristics can remain strong in easier regimes, whereas uncertainty-gated control becomes more valuable as disturbance structure and switching pressure increase.

\subsection{Real-Trace Replay}

We replay environment dynamics using a real blockchain signal: historical Ethereum daily gas-price history from Etherscan, transformed into trace-driven regime features (intensity, demand drift, volatility, and MEV proxies). Reported values in this subsection are from a single fixed historical-trace replay per controller under a deterministic schedule (seed count $n{=}1$ for this trace report), so these numbers should be interpreted as case-study evidence rather than variance-estimated statistics. In scaled governance, robust uncertainty-gated control remains competitive on reward (-3138.95 vs -3108.62 heuristic and -3265.31 PPO-Lagrangian) and keeps lower latency than PPO-Lagrangian (1.103 vs 1.253). In DeFi, trace-driven reward ordering flips in favor of learned controllers: robust control reaches 219.79, PPO-Lagrangian reaches 206.38, and heuristic control reaches 177.19 under the same historical trace schedule.

\begin{table}[htbp]
\centering
\input{tables/cost_of_safety_summary}
\caption{Explicit cost-of-safety quantification under matched protocol settings.}
\label{tab:cost-of-safety}
\end{table}

\subsection{Uncertainty Head-to-Head}

\begin{table}[htbp]
\centering
\input{tables/uncertainty_head_to_head}
\caption{Uncertainty method comparison at matched 90th-percentile trigger rates (10\%).}
\label{tab:uncertainty-head-to-head}
\end{table}

Using matched 90th-percentile trigger rates (10\% for every method), quantum-inspired parameter-perturbation uncertainty improves precision over classical expensive estimators at the same trigger budget (0.942 vs 0.833 for MC-dropout and 0.814 for ensemble), while preserving similar recall (0.106 vs 0.094 and 0.092). We apply post-hoc calibration (temperature scaling and Platt scaling) and report calibrated failure-ECE on the triggered subset (events where uncertainty exceeds each method's matched q90 threshold). Calibrated failure-ECE is below 0.1 for all methods in this protocol (values shown to four decimals in Table~\ref{tab:uncertainty-head-to-head}). In the hardest regime (intensity 1.0), quantum-inspired perturbation reaches 0.964 precision at 0.101 recall, while MC-dropout reaches 0.743 precision at 0.100 recall.
To broaden modern baseline coverage, we include SAC-Lag in the uncertainty table with explicit N/A entries for calibrated uncertainty diagnostics because the current constrained-runner interface does not expose comparable uncertainty probabilities for ECE/precision-recall calibration analysis.

\begin{figure}[H]
\centering
\includegraphics[width=0.82\linewidth]{figures/uncertainty_reliability.pdf}
\caption{Reliability visualization for calibrated uncertainty-triggered failure prediction.}
\label{fig:uncertainty-reliability}
\end{figure}

Figure~\ref{fig:uncertainty-reliability} complements Table~\ref{tab:uncertainty-head-to-head} by showing that calibration remains well-behaved across bins after post-hoc scaling.

\subsection{Safety-Reward Pareto Frontier}

\begin{figure}[H]
\centering
\includegraphics[width=0.82\linewidth]{figures/safety_reward_pareto.pdf}
\caption{Safety-reward Pareto view (scaled environment): robust gated controller vs CPO and P3O.}
\label{fig:safety-reward-pareto}
\end{figure}

Figure~\ref{fig:safety-reward-pareto} provides a direct trade-off visualization alongside Table~\ref{tab:constrained-baselines}. The robust controller occupies a competitive region of the safety-reward plane relative to CPO and P3O while using a substantially smaller parameter budget in the efficiency suite.

\subsection{Parameter Efficiency}

\begin{figure}[H]
\centering
\includegraphics[width=0.82\linewidth]{figures/parameter_efficiency_pareto.pdf}
\caption{Reward versus parameter count (log-scale) for classical MLP and VQC policy families.}
\label{fig:pareto-top-tier}
\end{figure}

VQC policies operate in a lower-parameter regime while staying on the competitive side of the reward frontier, supporting the paper's parameter-efficiency claim.

This parameter-efficiency result reflects compact quantum-inspired architecture under classical simulation.

The larger-model efficiency comparison (the protocol used for the 166.89$\times$ claim) quantifies this gap directly: classical MLP uses 4673 parameters while VQC (6-qubit, 4-layer) uses 28 (166.89$\times$ fewer). In that protocol, VQC improves mean reward (-2700.23 vs -2802.85) with lower variance (34.78 vs 123.20) at effectively identical safety violations (115.50 vs 115.50).
Expressed as relative context for that larger-model protocol, the VQC reward is 3.66\% better than the classical MLP baseline while using 0.60\% of the parameters.

\subsection{Matched-Capacity Efficiency Control}

\begin{table}[htbp]
\centering
\input{tables/parameter_efficiency_matched}
\caption{Matched-capacity efficiency control on scaled governance (VQC-28 vs MLP-28, plus larger MLP reference). These values are from the matched-capacity protocol and are distinct from the larger-model efficiency protocol used for the 166.89$\times$ compression claim.}
\label{tab:matched-capacity}
\end{table}

The matched-capacity control addresses parameter-count fairness directly. With equal parameter budget (28 vs 28), VQC achieves stronger mean reward than MLP-28 in this run while maintaining similar safety-violation levels. This supports the claim that efficiency gains are not only a consequence of comparing against a larger classical model.
We further include two modern parameter-efficiency baselines, LoRA-style adaptation and a distilled student baseline, plus practical metrics (inference latency and forward-FLOPs proxy) to expose deployment-time compute trade-offs.
The deployment trade-off is explicit in Table~\ref{tab:matched-capacity}: VQC inference latency is substantially higher than MLP-28 in this implementation (9.3438 ms vs 0.0837 ms). This is not ignored by the method; it defines the target regime. The intended use case is footprint-constrained or uncertainty-sensitive control loops where parameter budget and uncertainty quality matter more than raw per-step latency. For strict low-latency settings, compact classical baselines remain the preferred deployment option.

\subsection{Architecture View}

\begin{figure}[H]
\centering
\includegraphics[width=0.95\linewidth]{figures/architecture_loop_top_tier.pdf}
\caption{End-to-end governance loop: policy proposal, quantum-inspired uncertainty gate, shield enforcement, and on-chain commit.}
\label{fig:arch-top-tier}
\end{figure}

\FloatBarrier